\documentclass[a4paper,12pt]{article}
\usepackage[T2A]{fontenc}
\usepackage[utf8]{inputenc}
\usepackage[russian]{babel}
\usepackage{amsmath,amssymb,mathtools}
\usepackage{icomma}
\usepackage{geometry}
\usepackage{microtype}
\usepackage{tikz}
\usetikzlibrary{arrows.meta,positioning,shapes.geometric,calc}
\usepackage{booktabs,tabularx,array}
\usepackage{enumitem}
\usepackage{fancyhdr}
\usepackage{setspace}
\usepackage{xcolor}
\usepackage{needspace}
\usepackage[hidelinks]{hyperref}
\geometry{top=21mm,bottom=23mm,left=23mm,right=23mm,headheight=24pt,headsep=10pt}
\definecolor{accent}{HTML}{007E84}
\definecolor{darkaccent}{HTML}{005C63}
\definecolor{bodytext}{HTML}{272B2B}
\definecolor{lightgray}{HTML}{697376}
\color{bodytext}
\setlength{\parindent}{0pt}
\setlength{\parskip}{5pt}
\setstretch{1.24}
\setlist[itemize]{leftmargin=7mm,itemsep=2pt,topsep=2pt}
\setlist[enumerate]{leftmargin=9mm,itemsep=2pt,topsep=3pt}
\fancypagestyle{main}{%
\fancyhf{}%
\fancyhead[L]{\footnotesize\color{lightgray} Дроби и порядок действий}%
\fancyhead[R]{\footnotesize\color{lightgray} Надя Клименко}%
\fancyfoot[C]{\small\thepage}%
\renewcommand{\headrulewidth}{0.35pt}%
\renewcommand{\headrule}{\hbox to\headwidth{\color{accent}\leaders\hrule height \headrulewidth\hfill}}%
}
\pagestyle{main}
\newcommand{\sectionhead}[1]{%
  \needspace{4\baselineskip}
  \vspace{0.9\baselineskip}
  {\large\color{darkaccent}#1\par}
  \vspace{2pt}{\color{accent}\hrule height .5pt}\vspace{6pt}
}
\newcommand{\minorhead}[1]{\vspace{4pt}{\color{darkaccent}\bfseries #1}\par\vspace{-2pt}}
\newcommand{\checkitem}[1]{\item[$\square$] #1}
\newcommand{\result}[1]{\textcolor{darkaccent}{\textbf{#1}}}
\newcommand{\dec}[1]{#1}

\begin{document}
\thispagestyle{empty}
\begin{center}
\vspace*{28mm}
{\large\color{darkaccent} МАТЕМАТИКА}\par
\vspace{22mm}
{\LARGE Чек-лист}\par
\vspace{15mm}
{\LARGE\color{accent} Сложные и десятичные дроби}\\[4pt]
{\LARGE\color{accent} Порядок действий и сокращение}\par
\vspace{17mm}
{\normalsize Повторение методов, разобранных на занятии}\par
\vspace{29mm}
{\large Надя Клименко}\par
\vspace{13mm}
{\normalsize 10 ноября 2026 г.}\par
\vspace{22mm}
{\small Лёвин Артём Александрович}\\[2pt]
{\small эксклюзивно для Нади Клименко}\par
\vfill
\begin{tikzpicture}[x=1cm,y=1cm]
\draw[accent,line width=.8pt] (0,0)
 .. controls (1.3,.47) and (2.2,-.46) .. (3.35,0)
 .. controls (4.6,.52) and (5.2,-.50) .. (6.5,0)
 .. controls (7.7,.47) and (8.4,-.39) .. (9.6,0);
\fill[accent] (4.83,-.115) circle (1pt);
\end{tikzpicture}
\vspace{6mm}
\end{center}
\clearpage

\sectionhead{1. Что нужно уметь после занятия}
\begin{itemize}[label={},leftmargin=0mm,itemsep=3pt]
\checkitem{Преобразовывать сложную (\textit{многоэтажную}) дробь тремя способами.}
\checkitem{Переводить смешанные и десятичные числа в обыкновенные дроби.}
\checkitem{Убирать десятичную запятую из знаменателя без изменения значения дроби.}
\checkitem{Соблюдать порядок деления и умножения, учитывать скобки.}
\checkitem{Сокращать общие множители до громоздких вычислений.}
\checkitem{При необходимости раскладывать большие числа на простые множители.}
\end{itemize}

\sectionhead{2. Сложная дробь: три способа}
\textbf{Сложная дробь} --- дробь, в числителе или знаменателе которой также есть дробь. Основное условие: делитель и каждый знаменатель должны быть отличны от нуля.

\[\frac{\frac{a}{b}}{\frac{c}{d}}=\frac{a}{b}:\frac{c}{d}=\frac{ad}{bc},
\qquad b\ne0,\ c\ne0,\ d\ne0.\]

Рассмотрим один пример всеми способами:
\[\frac{1\frac35}{\frac23}=\frac{\frac85}{\frac23}.\]

\minorhead{Способ 1. Умножение числителя и знаменателя на одно число}
Выбираем число, кратное знаменателям $5$ и $3$: это $15$.
\[
\frac{\frac85}{\frac23}
=\frac{\frac85\cdot15}{\frac23\cdot15}
=\frac{24}{10}=\frac{12}{5}=2\frac25=2{,}4.
\]
Множитель $15$ действует \textbf{на весь числитель и на весь знаменатель}.

\minorhead{Способ 2. Перекрёстное умножение}
\[
\frac{\frac85}{\frac23}=\frac{8\cdot3}{5\cdot2}
=\frac{24}{10}=2{,}4.
\]
Этот способ особенно удобен, когда в дробях появляются буквы и скобки.

\minorhead{Способ 3. Деление через обратную дробь}
\[
\frac85:\frac23=\frac85\cdot\frac32
=\frac{24}{10}=2{,}4.
\]
\result{Итог:} все три приёма дают один ответ. Для обычных чисел выбирайте наиболее удобный.

\sectionhead{3. Трёхэтажная дробь и «невидимая единица»}
Целое число всегда можно представить дробью со знаменателем $1$:
\[
\frac{\frac72}{3}=\frac{\frac72}{\frac31}=\frac7{2\cdot3}=\frac76,
\qquad
\frac7{\frac23}=\frac{\frac71}{\frac23}=\frac{7\cdot3}{2}=\frac{21}{2}.
\]
Проверяйте, \textbf{где именно} находится дробь: в числителе или в знаменателе.

\sectionhead{4. Смешанные и десятичные дроби}
\minorhead{Смешанная дробь $\longrightarrow$ неправильная}
\[
2\frac57=\frac{2\cdot7+5}{7}=\frac{19}{7},
\qquad
2\frac13=\frac{2\cdot3+1}{3}=\frac73.
\]
Числитель равен произведению целой части и знаменателя \textbf{плюс} числитель дробной части.

\minorhead{Десятичная дробь $\longrightarrow$ обыкновенная}
\[
0{,}25=\frac{25}{100}=\frac14,\qquad
1{,}2=\frac{12}{10}=\frac65,\qquad
0{,}005=\frac5{1000}.
\]
Для деления удобно перенести запятую \textbf{в обоих числах} на одинаковое количество знаков:
\[
\frac{3{,}6}{0{,}12}=\frac{360}{12}=30,
\qquad
\frac{0{,}020}{0{,}005}=\frac{20}{5}=4.
\]
В первом примере умножили числитель и знаменатель на $100$, во втором --- на $1000$.

\sectionhead{5. Скобки меняют делитель}
Умножение и деление без скобок выполняют \textbf{слева направо}:
\[
a:b\cdot c=\frac{a}{b}\cdot c=\frac{ac}{b},
\qquad
 a:(b\cdot c)=\frac{a}{bc},\qquad b,c\ne0.\]

\textbf{Сравните результаты:}
\[
12:3\cdot2=4\cdot2=8,
\qquad
12:(3\cdot2)=12:6=2.
\]
При переносе множителя в числитель убедитесь, что он действительно стоит \textbf{после деления}, вне скобок делителя.

\sectionhead{6. Сокращение до умножения}
В выражении из занятия смешанная дробь сначала превращается в неправильную:
\[
12{,}3:2{,}8\cdot2\frac13
=\frac{12{,}3}{2{,}8}\cdot\frac73
=\frac{123\cdot7}{28\cdot3}
=\frac{41}{4}=10{,}25.
\]
Сократили $123$ и $3$, затем $7$ и $28$. Так вычисления стали короче.

\textbf{Правило:} сокращают \textbf{множители} числителя и знаменателя. В сумме или разности сначала требуется преобразование, например вынесение общего множителя:
\[
\frac{2+6}{2}=4,
\qquad \text{нельзя получить }1+6=7\text{ «сокращением двойки».}
\]

\sectionhead{7. Сумма дробей с десятичным знаменателем}
Вначале избавьтесь от десятичных знаменателей, потом сократите и сложите:
\begin{align*}
\frac2{0{,}9}+\frac5{1{,}8}
&=\frac{20}{9}+\frac{50}{18}\\
&=\frac{20}{9}+\frac{25}{9}
=\frac{45}{9}=5.
\end{align*}
Выбирайте общий знаменатель \textbf{после сокращения}: здесь он сразу стал равен $9$.

\sectionhead{8. Много множителей: разложение на простые}
Если большие числа сокращаются с трудом, замените их произведениями простых множителей.
\[
\frac{0{,}16\cdot1{,}44\cdot0{,}42}{0{,}72\cdot0{,}32\cdot3{,}78}
=\frac{16\cdot144\cdot42}{72\cdot32\cdot378}.
\]
Здесь в каждом произведении три множителя с двумя знаками после запятой; значит, \textbf{весь} числитель и \textbf{весь} знаменатель умножены на $100^3$.

\begin{align*}
16&=2^4, &144&=2^4\cdot3^2, &42&=2\cdot3\cdot7,\\
72&=2^3\cdot3^2, &32&=2^5, &378&=2\cdot3^3\cdot7.
\end{align*}
Тогда
\[
\frac{2^4\cdot 2^4\cdot3^2\cdot2\cdot3\cdot7}
{2^3\cdot3^2\cdot2^5\cdot2\cdot3^3\cdot7}
=\frac{2^9\cdot3^3\cdot7}{2^9\cdot3^5\cdot7}
=\boxed{\frac19}.
\]
\textbf{Когда разлагать?} Если чисел мало и сокращение очевидно --- сокращайте сразу. Если множителей много --- разложение снижает риск ошибки.

\sectionhead{9. Комбинированный пример: действуйте по частям}
Если многоэтажная дробь содержит разные действия, сначала вычислите её числитель и знаменатель отдельно:
\[
\frac{\displaystyle 1\frac1{10}\cdot\frac76:\frac7{60}}
{\displaystyle \frac17:0{,}25\cdot2\frac25}.
\]
\textbf{Числитель:} переводим смешанное число, затем сокращаем.
\[
1\frac1{10}\cdot\frac76:\frac7{60}
=\frac{11}{10}\cdot\frac76\cdot\frac{60}{7}=11.
\]
\textbf{Знаменатель:} деление и умножение выполняются слева направо.
\[
\frac17:0{,}25\cdot2\frac25
=\frac17\cdot4\cdot\frac{12}{5}=\frac{48}{35}.
\]
\textbf{Последний шаг:} делим число на дробь.
\[
11:\frac{48}{35}=11\cdot\frac{35}{48}
=\boxed{\frac{385}{48}=8\frac1{48}}.
\]

\sectionhead{10. Типичные ошибки и самопроверка}
\begin{itemize}[label=\textcolor{accent}{$\blacktriangleright$},leftmargin=6mm]
\item \textbf{Неравное домножение:} $\frac ab\ne\frac{100a}{10b}$ в общем случае. Числитель и знаменатель домножают на \textbf{одно и то же} ненулевое число.
\item \textbf{Ошибка в смешанной дроби:} $2\frac57=\frac{19}{7}$; число $2\cdot7+5$ считайте до сокращений.
\item \textbf{Потеря скобок:} $a:b\cdot c$ и $a:(b\cdot c)$ обычно различаются.
\item \textbf{Сокращение суммы:} зачёркивать одинаковые части слагаемых как множители нельзя.
\item \textbf{Деление на нуль:} знаменатели и делители должны быть отличны от нуля.
\item \textbf{Проверка результата:} оцените порядок величины и подставьте ответ в исходное выражение.
\end{itemize}
\vspace{2pt}
\textcolor{darkaccent}{\textbf{Небольшое повторение процентов.}}
$20\%$ от $500$ --- это $500\cdot\frac{20}{100}=100$.

\clearpage
\begin{center}
{\Large\color{darkaccent} Алгоритм: как вычислить сложное дробное выражение}\par
\vspace{10mm}
\end{center}
\begin{center}
\begin{tikzpicture}[
  font=\small,
  >=Stealth,
  node distance=8mm,
  startstop/.style={ellipse,draw=accent,thick,align=center,minimum width=46mm,minimum height=9mm,inner sep=2mm},
  stagebox/.style={rectangle,rounded corners=2mm,draw=accent,align=center,text width=80mm,minimum height=10mm,inner sep=2.2mm},
  decision/.style={diamond,draw=darkaccent,aspect=3.1,align=center,text width=48mm,inner sep=1.4mm},
  side/.style={rectangle,rounded corners=2mm,draw=accent,align=center,text width=43mm,minimum height=10mm,inner sep=2mm},
  arr/.style={-Stealth,thick,color=darkaccent}
]
\node[startstop] (s) {Начало};
\node[stagebox,below=of s] (order) {Проверить скобки и порядок действий.\\Определить границы делителя.};
\node[decision,below=9mm of order] (dec) {Есть десятичные\\знаменатели?};
\node[side,right=22mm of dec] (times) {Одинаково домножить верх и низ дроби на $10^n$};
\node[stagebox,below=11mm of dec] (mix) {Перевести смешанные дроби\\в неправильные; при необходимости\\десятичные --- в обыкновенные.};
\node[stagebox,below=of mix] (rev) {Заменить деление умножением\\на обратную дробь.};
\node[decision,below=9mm of rev] (heavy) {Множителей много,\\числа велики?};
\node[side,right=22mm of heavy] (factor) {Разложить числа на простые множители};
\node[stagebox,below=11mm of heavy] (cancel) {Сократить общие множители\\в числителе и знаменателе.};
\node[stagebox,below=of cancel] (finishstep) {Вычислить оставшееся, проверить знаки, скобки и величину ответа.};
\node[startstop,below=of finishstep] (f) {Ответ};
\draw[arr] (s) -- (order);
\draw[arr] (order) -- (dec);
\draw[arr] (dec) -- node[left,font=\scriptsize]{нет} (mix);
\draw[arr] (dec) -- node[above,font=\scriptsize]{да} (times);
\draw[arr] (times.south) |- ([yshift=4mm]mix.north) -- (mix.north);
\draw[arr] (mix) -- (rev);
\draw[arr] (rev) -- (heavy);
\draw[arr] (heavy) -- node[left,font=\scriptsize]{нет} (cancel);
\draw[arr] (heavy) -- node[above,font=\scriptsize]{да} (factor);
\draw[arr] (factor.south) |- ([yshift=4mm]cancel.north) -- (cancel.north);
\draw[arr] (cancel) -- (finishstep);
\draw[arr] (finishstep) -- (f);
\end{tikzpicture}
\end{center}
\clearpage

\sectionhead{Тренировочный блок --- 10 заданий}
\textcolor{lightgray}{Решите самостоятельно. Сначала определите удобный способ преобразования; в последних заданиях покажите ключевые промежуточные действия.}
\par\vspace{5mm}
\textbf{Базовый уровень}\par
\begin{enumerate}[label=\textcolor{accent}{\arabic*.},itemsep=9pt]
\item Вычислите: $\displaystyle\frac{1\frac12}{\frac34}$.
\item Вычислите: $\displaystyle\frac7{\frac23}$.
\item Упростите: $\displaystyle\frac{\frac56}{\frac{10}{9}}$.
\item Вычислите: $\displaystyle\frac{0{,}020}{0{,}005}$.
\item Вычислите: $\displaystyle3{,}6:0{,}12$.
\end{enumerate}
\vspace{4mm}
\textbf{Повышенный уровень}\par
\begin{enumerate}[label=\textcolor{accent}{\arabic*.},start=6,itemsep=10pt]
\item Вычислите: $\displaystyle\frac4{0{,}8}+\frac3{1{,}2}$.
\item Вычислите: $\displaystyle7{,}2:2{,}4\cdot1\frac13$.
\item Сравните с предыдущим заданием и вычислите: $\displaystyle7{,}2:\bigl(2{,}4\cdot1\frac13\bigr)$.
\end{enumerate}
\vspace{4mm}
\textbf{Сложный уровень}\par
\begin{enumerate}[label=\textcolor{accent}{\arabic*.},start=9,itemsep=12pt]
\item Вычислите, используя сокращение или разложение на простые множители:
\[\frac{0{,}16\cdot1{,}44\cdot0{,}42}{0{,}72\cdot0{,}32\cdot3{,}78}.\]
\item Вычислите по действиям (сначала числитель и знаменатель):
\[\frac{\displaystyle 1\frac15\cdot\frac79:\frac7{30}}
{\displaystyle \frac18:0{,}25\cdot1\frac34}.\]
\end{enumerate}
\vfill
{\small\color{lightgray} Проверьте: оба числа при домножении меняются одинаково; порядок действий соответствует скобкам; сокращались только множители.}
\clearpage

\sectionhead{Ответы к тренировочному блоку}
\begin{center}
\renewcommand{\arraystretch}{1.55}
\begin{tabular}{@{}>{\centering\arraybackslash}p{22mm}p{94mm}@{}}
\toprule
\textbf{№} & \textbf{Ответ} \\
\midrule
1 & $2$ \\
2 & $\dfrac{21}{2}=10\dfrac12$ \\
3 & $\dfrac34$ \\
4 & $4$ \\
5 & $30$ \\
6 & $7{,}5$ \\
7 & $4$ \\
8 & $2{,}25$ \\
9 & $\dfrac19$ \\
10 & $\dfrac{32}{7}=4\dfrac47$ \\
\bottomrule
\end{tabular}
\end{center}
\vspace{13mm}
\minorhead{Итоговая памятка}
\begin{tabularx}{\linewidth}{@{}>{\raggedright\arraybackslash\color{darkaccent}}p{47mm}>{\raggedright\arraybackslash}X@{}}
\toprule
\multicolumn{1}{@{}l}{\textbf{Сигнал в выражении}} & \multicolumn{1}{l@{}}{\textbf{Что делать}}\\
\midrule
Дробь внутри дроби & Замените большую черту делением или примените перекрёстное умножение.\\
Смешанное число & Переведите в неправильную дробь.\\
Запятая в знаменателе & Домножьте числитель и знаменатель на одно и то же $10^n$.\\
Умножение рядом с делением & Проверьте скобки и выполняйте действия слева направо.\\
Много множителей & Сократите общие множители; при необходимости разложите на простые.\\
Сумма дробей & Сначала сократите отдельные дроби, потом найдите общий знаменатель.\\
\bottomrule
\end{tabularx}
\vfill
\begin{center}
{\small\color{lightgray} Лёвин Артём Александрович --- эксклюзивно для Нади Клименко}\par
\end{center}
\end{document}
